% F.1 companion ledger, copied from the supplied transcription files.
% Scan numbers, rather than the sometimes shifted manuscript foliation, are used.
\subsection{Manuscript statements settled by Parts D and E}

\begin{remark}[D.1, cote 8, scan 81]\label{rem:DE-quote-D1}
The page says ``Elle est évidemment nulle si $M_0$ est de rang $\leq1$, mais sans doute pas nulle si $M_0$ de rang $d\geq2$,'' in the universal general-linear case.  Theorem~\ref{thm:DE-split} confirms the obstruction interpretation and corrects this proposed nonsplitting list by the three exceptions $(d,p)=(2,2),(2,3),(3,2)$.  It also gives the requested complete special-linear and symplectic lists.  The distinction between a supplement and a complement is retained.
\end{remark}

\begin{remark}[D.2, cote 8, scans 80--81]\label{rem:DE-quote-D2}
At scan 81 the transcribed comparison ends in $H^2(S,\operatorname{End}(E))$, while the concluding argument says ``si le rev. $S'$ de $S$ est le rev. universel'' then the group-cohomology edge map ``est injectif.''  These are distinct steps.  Theorem~\ref{thm:DE-coherent-obstruction} constructs the requested smooth projective $S$ with $\xi\ne0$ for every nonsplit pair.  Remark~\ref{rem:DE-universal-negative} constructs full-monodromy smooth projective curves with $\xi=0$, so the universal version of Lead D is false.
\end{remark}

\begin{remark}[E.1 and R11, cote 5, scan 41]\label{rem:DE-quote-E1}
After the two rules for $\pi$, the question is ``Si on a une application satisfaisant ces conditions, peut on définir les $\pi^{(n)}$'' satisfying the divided-power axioms.  Theorem~\ref{thm:DE-PD-pi} answers yes for every $\mathbb Z_{(p)}$-algebra, including arbitrary $p$-torsion and $p=2$.  The same page asks of $c_n\bmod p$, ``Est-ce tjs $1$?''  The answer is no for $p\geq3$: the exact formula gives $c_n\equiv\prod_i(a_i!)^{-1}\pmod p$ for the base-$p$ digits $a_i$, and already $n=2$ gives $1/2\ne1$.  At $p=2$ all these residues are $1$.
\end{remark}

\begin{remark}[E.2, cote 6, scan 7, section 1.8]\label{rem:DE-quote-E2}
The exact transcribed sentence reads ``Sauf erreur, on a encore une équivalence de catégories si $S\to R$ est plat et localement de présentation finie,'' followed by the insertion ``(du moins si R loc.noeth., et se bornant aux Modules cohérents)'', then ``mais je n'ai pas écrit'' the demonstration.  The insertion materially narrows the manuscript conjecture.  Proposition~\ref{prop:DE-E2-coherent} proves the coherent characteristic-zero field case.  Proposition~\ref{prop:DE-E2-Fpure} proves the arbitrary quasi-coherent case for $F$-pure schemes locally of finite type over a perfect field of positive characteristic, using Lemma~\ref{lem:DE-crystal-Fdiv}.  The full flat relative assertion requested in Part E.2 is refuted by Theorem~\ref{thm:DE-E2-counterexample}, which gives a non-Noetherian finite free example and a rank-one crystal automorphism invisible to restriction.  Corollary~\ref{cor:DE-E2-Noetherian-counterexample} further gives nonfaithfulness over a regular Noetherian base with arbitrary quasi-coherent coefficients.  The inserted relative coherent conjecture over general locally Noetherian bases remains the distinct undecided refinement.
\end{remark}

\begin{remark}[E.3 and R10, cote 7, scan 17]\label{rem:DE-quote-E3}
The hoped-for comparison with integral de Rham cohomology is followed by ``C'est faux malheureusement'' and the ordinary-abelian calculation
$H^i((X_0/S)_{\mathrm{cris}},\underline O)=0$ for $i>\dim X_0$.  Here the notation $\mathrm{cris}$ means the old infinitesimal site without divided powers.  Theorem~\ref{thm:DE-inf} identifies this cohomology integrally with $H^i_{\et}(X_0,\Zp)\otimes_{\Zp}W$, including torsion.  Proposition~\ref{prop:DE-R10} verifies the displayed vanishing and the preceding mod-$p$ statement.  No identification with modern PD crystalline cohomology is intended.
\end{remark}

\begin{remark}[E.4, cote 7, scan 70]\label{rem:DE-quote-E4}
The conjecture says ``Pour tout espace analytique $\mathfrak X$, le complexe $\mathcal F_{\mathfrak X}^*$ est une résolution du faisceau $\mathbf C_{\mathfrak X}$.''  Theorem~\ref{thm:DE-intrinsic-nerve} settles the algebraic finite-type characteristic-zero question posed in Part E.4: the intrinsic completed nerve computes infinitesimal de Rham cohomology, including for singular schemes.  For $\mathbb C[t^2,t^3]$, it gives $H^0=\mathbb C$ and $H^i=0$ for $i>0$.  The theorem is not stated as a proof of the broader assertion for every analytic space.
\end{remark}

\begin{remark}[E.5, cote 7, scan 24]\label{rem:DE-quote-E5}
For a singular complex curve, the page says ``on n'obtient déjà pas pour'' $\dim H^1_{\mathrm{DR}}(X)$ ``le premier nb de Betti comme il faudrait.''  Its proposed ordinary-cusp example does not establish this.  Theorem~\ref{thm:DE-curves} proves the precise correction for reduced affine locally planar complex curves:
\[
 \dim_{\mathbb C}H^1(\Omega_C^\bullet)
 =b_1(C)+\sum_x(\mu_x-\tau_x),\qquad
 H^i(\Omega_C^\bullet)=0\quad(i\geq2).
\]
Thus the extra term vanishes for the ordinary cusp, although it need not vanish for other plane singularities.
\end{remark}

\begin{remark}[E.6, cote 7, scans 42--43]\label{rem:DE-quote-E6}
The proposed theorem begins ``Il semble que le théorème suivant soit vrai.''  The key assertion at scan 43 is ``Cet homomorphisme ne dépend que de l'homomorphisme $u_0$'' induced modulo $p$; the page notes that the then-available proof was only ``mod torsion, i.e. en tensorisant par $K$.''  Theorem~\ref{thm:DE-integral-MW-independence} confirms the integral assertion, for every prime and in every degree, with the required overconvergent homotopy supplied by Ertl--Sprang.
\end{remark}

\begin{remark}[E.7, cote 8, scans 48--50]\label{rem:DE-quote-E7}
The notation is $t_G=\pi^{-\alpha}\Omega^\vee$.  At scan 50, after comparing the two full lattices, the page asks ``On peut se demander si ces sous-réseaux sont les mêmes'' and adds ``ce qui impliquerait évidemment'' $\alpha=0$.  Proposition~\ref{prop:DE-LT} proves $\alpha=0$ for every finite extension $K/\Qp$, but refutes equality of the full lattices using $\pi^{p+1}=p$.  The implication on the manuscript page is correct; a converse or equivalence must not be attributed to it.  The dual $\Omega^\vee$ is the $\Zp$-dual, so the inverse different is part of the normalization.
\end{remark}

\begin{remark}[E.8, cote 9, scan 22]\label{rem:DE-quote-E8}
The assertion is ``cet isom est compatible avec les $p^{-1}$-structures, donc transforme $\sigma_A$ en $\sigma_{A^*}$.''  Theorem~\ref{thm:DE-hasse} proves exact transport of the Hasse invariant under the stated Hodge-determinant and trace identification, for every abelian scheme over an $\Fp$-scheme, with no reducedness or ordinarity assumption.
\end{remark}

\begin{remark}[E.9, cote 9, scan 29]\label{rem:DE-quote-E9}
Part c) reads ``Il faudrait vérifier que'' $\alpha_{A^*}=\alpha_A^{-1}$, followed by ``[ aux signes (?)\ldots].''  Part d) compares the reduction of $\alpha_A$ with the map ``déduit plus haut de la dualité en cohomologie de DeRham (à vérifier \ldots).''  Theorem~\ref{thm:DE-ordinary-sign} proves exact inverse biduality for standard Poincaré biduality, and computes the de Rham comparison factor
$\varepsilon_g=(-1)^{g(g+1)/2}$ with its required direction/inverse.  In particular the literal lead's simple multiplicativity is false in odd characteristic: a product contributes the additional factor $(-1)^{gh}$.
\end{remark}

\begin{remark}[E.10, cote 7, scans 18--19]\label{rem:DE-quote-E10}
The heading is ``Tests décisifs'': for smooth affine $X_0$, ``on doit retrouver la définition de W.M.''; for smooth proper liftable $X_0$, ``on doit retrouver (du moins mod torsion) la cohomologie de De Rham''.  The page calls this ``une première'' ``esquisse.''  Proposition~\ref{prop:DE-old-site}, Theorem~\ref{thm:DE-old-Gm}, and Proposition~\ref{prop:DE-old-proper} refute the integral affine test, its rational version, and the proper test, respectively, on the literal flat weak-formal site.  The rational rigid-cohomology claim of Lead E10 is therefore false as well.
\end{remark}

\begin{remark}[E.11, cote 6, scan 14]\label{rem:DE-quote-E11}
The page concludes ``On voit donc là un principe d'approche transcendante pour l'étude des familles de groupes formels'' in characteristic $p$.  Theorem~\ref{thm:DE-simply-connected-convergent} proves the requested statement for convergent $F$-isocrystals under the specified proper smooth algebraic lift and simply connected complex fibre.  Corollary~\ref{cor:DE-BT-constant} proves the additional isogeny-to-constant assertion for Barsotti--Tate groups in the literal lead.  These hypotheses retain both the specified embedding $W(k)\hookrightarrow\mathbb C$ and the condition on the generic complex fibre.
\end{remark}
